\documentclass[11pt]{article}

\usepackage[a4paper,margin=1in]{geometry}
\usepackage{amsmath,amssymb,amsthm}
\usepackage{graphicx}
\usepackage{booktabs}
\usepackage{algorithm}
\usepackage{algpseudocode}
\usepackage{hyperref}

\title{Greedy Walks on Neighbour-Hyperplane Graphs\\in $p$-adic Regression (and a Sudoku Analogue)}
\author{(Author details omitted)}
\date{\today}

\newtheorem{definition}{Definition}
\newtheorem{lemma}{Lemma}

\begin{document}
\maketitle

\begin{abstract}
In ultrametric regression problems (including $p$-adic regression), the objective is typically locally constant and ``gradient-like'' training rules are uninformative. One natural alternative is a discrete greedy descent on a finite set of candidate models.
This note records a simple computational experiment motivated by the following phenomenon: when the loss is ultrametric, a best-fit hyperplane can be chosen to pass through $d{+}1$ data points (for $d$ features), so there are finitely many candidates. If we define two hyperplanes to be \emph{neighbours} when they differ by exactly one of their defining points, we obtain a natural graph on candidate hyperplanes. A naive optimiser then repeatedly moves to the neighbouring hyperplane with the largest loss decrease, stopping at a local minimum.

We implement this neighbour-hyperplane graph for a toy $p$-adic linear regression model with synthetic $p$-adic noise, enumerate all hyperplanes through $d{+}1$ points, \emph{deduplicate identical coefficient vectors}, and measure how often several \emph{improvement-only descent policies} reach a \emph{global} minimum from a random start. Besides deterministic steepest descent, we consider randomised policies that choose among improving moves uniformly, proportional to improvement size, or via a softmax with temperature. We also include a noiseless baseline, where all points lie exactly on the generating hyperplane.
We also run an analogue experiment for a Sudoku solver whose move set is within-row swaps: the induced swap graph plays the role of a neighbour graph, and greedy descent corresponds to always taking the best improving swap until a local minimum.

Across the tested regimes, greedy descent frequently gets trapped: in the $p$-adic hyperplane graphs we observe many local minima and only moderate ``global basin'' fractions, and in Sudoku the pure greedy walk almost never solves harder carved instances unless multiple random restarts are allowed.
\end{abstract}

\section{Neighbour hyperplanes and greedy descent}
Fix a dataset of $n$ points $(x_i,y_i)$ with $x_i\in\mathbb{Z}^d$ and $y_i\in\mathbb{Z}$. A (affine) hyperplane is a parameter vector
\[
\beta = (\beta_0,\beta_1,\dots,\beta_d)
\]
representing the prediction $\hat y = \beta_0 + \sum_{j=1}^d \beta_j x_j$.

\begin{definition}[Ultrametric residual loss]
Fix a prime $p$ and define the $p$-adic absolute value on rationals by $|q|_p = p^{-v_p(q)}$ with $|0|_p=0$. For a parameter vector $\beta$ define the residuals
\[
r_i(\beta) = y_i - \beta_0 - \sum_{j=1}^d \beta_j x_{i,j},
\]
and the $\ell_1$ ultrametric loss
\[
L(\beta) = \sum_{i=1}^n |r_i(\beta)|_p.
\]
\end{definition}

The key property motivating a finite search is that, under ultrametric losses, best fits can be chosen to interpolate points. We take that as an organising principle: enumerate candidate hyperplanes by selecting $d{+}1$ points and solving the corresponding linear system over $\mathbb{Q}$.

\begin{definition}[Neighbour hyperplanes]
Let $S\subset\{1,\dots,n\}$ be a subset of indices of size $d{+}1$. Let $\beta(S)$ be the hyperplane interpolating the points in $S$ (if the system is nonsingular). Two such subsets $S$ and $S'$ are \emph{neighbours} if $|S\cap S'| = d$ (i.e.\ they differ by exactly one point).
\end{definition}

Different $(d{+}1)$-subsets can define the \emph{same} hyperplane (for example, in the noiseless setting where all points lie on one hyperplane). In our implementation we therefore \emph{deduplicate} by coefficient vector and treat the graph nodes as \emph{unique} hyperplanes; an edge between two hyperplanes exists if it is induced by at least one single-point swap between defining subsets.

Any \emph{improvement-only} descent policy repeatedly moves to a neighbouring node with strictly lower loss, stopping at a local minimum.

\begin{algorithm}[t]
\caption{Steepest-descent walk on neighbour hyperplanes}
\begin{algorithmic}[1]
\Require Candidate set $\{\beta(S)\}$ indexed by $(d{+}1)$-subsets $S$, loss $L(\cdot)$
\State Choose an initial subset $S$
\While{there exists a neighbour $S'$ of $S$ with $L(\beta(S')) < L(\beta(S))$}
  \State Set $S \gets \arg\min_{S' \sim S} L(\beta(S'))$ (break ties deterministically)
\EndWhile
\State \Return local minimum $S$
\end{algorithmic}
\end{algorithm}

\section{Synthetic $p$-adic regression experiments}
\subsection{Noise models}
There is no single canonical analogue of ``Gaussian noise'' on $\mathbb{Q}_p$. A standard baseline on $\mathbb{Z}_p$ is Haar-uniform noise. We use two simple integer-valued approximations:
\begin{itemize}
  \item \textbf{Haar (truncated)}: sample $U$ uniformly modulo $p^M$ and return $\varepsilon = p^{k_0}U$ (with a symmetric representative for $U$).
  \item \textbf{Valuation+unit}: sample a valuation $K\ge k_0$ with a geometric tail, choose a small coefficient $a$ not divisible by $p$, and return $\varepsilon = a\,p^K$.
\end{itemize}
Both reflect the idea that $v_p(\varepsilon)$ is typically $k_0$ but occasionally larger.

\subsection{Experimental design}
For each trial we sample an integer ``ground-truth'' hyperplane and integer covariates $x_i$, then set
\[
y_i = \beta_0^\star + \sum_{j=1}^d \beta_j^\star x_{i,j} + \varepsilon_i.
\]
In the runs reported here, the generating coefficients are sampled as $p$-powers (i.e.\ $\beta_j^\star \in \{\pm p^k\}$ for a small exponent range).

We then enumerate all hyperplanes through $d{+}1$ points, deduplicate identical coefficient vectors, compute $L(\beta)$ for each unique hyperplane, and analyse several improvement-only descent policies on the induced neighbour graph:
\begin{itemize}
  \item \textbf{Global-hit probability}: for a uniformly random start node, the probability that the descent process ends at a global minimum.
  \item \textbf{Local minima count}: the number of greedy sinks.
  \item \textbf{True-is-global}: whether the generating $\beta^\star$ is itself globally optimal under $L$.
\end{itemize}

The descent policies we test are:
\begin{itemize}
  \item \textbf{Steepest}: choose a neighbour with maximal loss decrease (deterministic tie-break).
  \item \textbf{Uniform}: choose uniformly among all improving neighbours.
  \item \textbf{Proportional}: choose with probability proportional to improvement size.
  \item \textbf{Softmax}: choose among improving neighbours with probability proportional to $\exp(\Delta/T)$ where $\Delta$ is the loss decrease and $T>0$ is a temperature (we use $T=1$).
\end{itemize}

In all runs we used $n=20$ points and $d=3$ features (so each candidate hyperplane is defined by $4$ points), coefficient and covariate bounds $5$, loss type $\ell_1$, and noise precision parameter $M=6$. For $p=3$ we ran $60$ trials per configuration; for $p=11$ we ran $20$. We additionally include a \textbf{noiseless} baseline ($\varepsilon_i\equiv 0$).

\subsection{Results}
Table~\ref{tab:padic-summary} summarises global-hit probabilities for each descent policy. Two qualitative findings stand out:
\begin{enumerate}
  \item With noise, the neighbour-hyperplane landscape has \emph{many} bad local minima; improvement-only descent reaches a global minimum only a fraction of the time (typically $\approx 0.15$--$0.40$ here), and the exact probability depends modestly on the policy.
  \item In the noisy settings tested here, the generating hyperplane $\beta^\star$ is never globally optimal. Intuitively, an interpolating hyperplane can ``spend'' degrees of freedom to make $d{+}1$ residuals exactly zero while keeping the remaining residual valuations high, so the optimum tends to overfit rather than recover $\beta^\star$. In contrast, in the noiseless baseline all hyperplanes coincide with $\beta^\star$ and every node is globally optimal.
\end{enumerate}

\begin{table}[t]
\centering
\begin{tabular}{@{}rrrcccc@{}}
\toprule
$p$ & model & $k_0$ & steepest & uniform & proportional & softmax ($T{=}1$) \\
\midrule
3  & none      & 0 & 1.000 & 1.000 & 1.000 & 1.000 \\
3  & haar      & 0 & 0.346 & 0.329 & 0.360 & 0.349 \\
3  & haar      & 1 & 0.346 & 0.331 & 0.361 & 0.339 \\
3  & haar      & 2 & 0.346 & 0.329 & 0.360 & 0.332 \\
3  & haar      & 3 & 0.346 & 0.329 & 0.360 & 0.330 \\
3  & valuation & 0 & 0.360 & 0.341 & 0.371 & 0.360 \\
3  & valuation & 1 & 0.360 & 0.341 & 0.371 & 0.349 \\
3  & valuation & 2 & 0.360 & 0.341 & 0.371 & 0.344 \\
3  & valuation & 3 & 0.360 & 0.341 & 0.371 & 0.342 \\
\midrule
11 & none      & 0 & 1.000 & 1.000 & 1.000 & 1.000 \\
11 & haar      & 0 & 0.221 & 0.189 & 0.210 & 0.213 \\
11 & haar      & 1 & 0.221 & 0.189 & 0.210 & 0.191 \\
11 & haar      & 2 & 0.221 & 0.189 & 0.210 & 0.189 \\
11 & haar      & 3 & 0.221 & 0.189 & 0.210 & 0.189 \\
11 & valuation & 0 & 0.227 & 0.192 & 0.211 & 0.214 \\
11 & valuation & 1 & 0.227 & 0.192 & 0.211 & 0.194 \\
11 & valuation & 2 & 0.227 & 0.192 & 0.211 & 0.192 \\
11 & valuation & 3 & 0.227 & 0.192 & 0.211 & 0.192 \\
\bottomrule
\end{tabular}
\caption{Global-hit probabilities for improvement-only descent policies on the neighbour-hyperplane graph ($n=20,d=3$). Each entry is the mean global-hit probability over trials for that configuration. (Standard errors and additional diagnostics are recorded in \texttt{outputs/2026-02-20/analysis/experiment\_report.md}.)}
\label{tab:padic-summary}
\end{table}

Figure~\ref{fig:padic-policy} highlights the noiseless-vs-noisy contrast and the modest policy dependence in a baseline noisy setting.

\begin{figure}[t]
\centering
\includegraphics[width=0.92\linewidth]{../outputs/2026-02-20/analysis/padic_lr_policy_noiseless_vs_noisy.pdf}
\caption{Noiseless vs baseline noisy ($\text{haar},k_0=0$) global-hit probabilities by policy. Error bars are standard errors across trials.}
\label{fig:padic-policy}
\end{figure}

\section{Sudoku as a neighbour-graph analogue}
The Sudoku solver in this repository is already phrased as a local-search problem on a discrete state space: initialise each row as a permutation of $\{1,\dots,9\}$ consistent with clues, and apply within-row swaps. This induces a natural neighbour graph (each state is a node; each swap is an edge) analogous to the neighbour-hyperplane graph above.

We consider two training dynamics:
\begin{itemize}
  \item \textbf{Pure greedy descent}: evaluate all swaps, take the best improving one, and stop at a local minimum.
  \item \textbf{Stepwise with escape}: take the best improving swap within a chosen row, but if no improving swap exists, apply a random swap to escape local minima.
\end{itemize}

\subsection{Experimental design}
We generated $200$ puzzles for each clue count $\{40,36,32,30,28,26,24\}$ by carving random solved grids (no uniqueness guarantee). For greedy descent we ran both one restart and $15$ restarts (each restart is a new random initialisation). For the stepwise-with-escape method we ran $50$ puzzles per clue count with a single restart and $20{,}000$ maximum steps.

\subsection{Results}
Figure~\ref{fig:sudoku-rate} summarises solve rates. The key qualitative conclusion is that \emph{pure greedy descent} is dominated by local minima: with one start it almost never solves, and even with $15$ restarts it solves only a minority of carved instances. By contrast, adding a simple escape mechanism makes solving routine in this carved setting.

\begin{figure}[t]
\centering
\includegraphics[width=0.80\linewidth]{../outputs/2026-02-20/analysis/sudoku_solve_rate.pdf}
\caption{Solve rates on carved Sudoku instances under three dynamics: greedy descent with one restart, greedy descent with $15$ restarts, and stepwise-with-escape.}
\label{fig:sudoku-rate}
\end{figure}

\section{Implementation notes}
All experiments are implemented as small, dependency-free Python programs:
\begin{itemize}
  \item \texttt{code/padic\_linear\_regression.py} generates data, enumerates hyperplanes, builds the neighbour relation, and computes greedy basins.
  \item \texttt{code/padic\_sudoku\_regression.py} contains Sudoku solvers, including \texttt{--method greedy} (pure descent) and \texttt{--method stepwise}.
  \item \texttt{code/run\_experiments.py} runs Sudoku trials and writes CSVs.
  \item \texttt{code/make\_experiment\_report.py} aggregates CSVs and produces the plots included here.
\end{itemize}

\section{Discussion}
The neighbour-hyperplane graph provides a concrete finite combinatorial object attached to an ultrametric regression instance, and greedy descent on this graph is an appealing ``first baseline'' heuristic. These experiments suggest two general cautions:
\begin{enumerate}
  \item \textbf{Local minima can be abundant.} Even when each node has only $O(nd)$ neighbours, the greedy-walk landscape can have many sinks.
  \item \textbf{Overfitting is structurally favoured.} When the candidate class contains interpolating models, ultrametric losses may strongly prefer exact interpolation of a subset of points unless an explicit regulariser counterbalances it.
\end{enumerate}
Both points appear again in Sudoku: if we remove escape moves, greedy descent is quickly trapped.

\end{document}
